% Isolated production-verification candidate; not live-book source.
% Proof of \cref{thm:p071-infinite-star-structure}.

\begin{proof}
We first prove \eqref{eq:p071-countable-full-branch-constancy}.  Define a
subsheaf \(\mathscr C\) of the sheaf of
\(\mathbf R\cup\{-\infty\}\)-valued functions on \(X_B\) as follows.  There
is no condition on an open set not containing \(\eta\).  If \(W\) contains
\(\eta\), a section \(v\) belongs to \(\mathscr C(W)\) when there is a
countable set \(C_v\subset B\) such that
\[
 W\cap X_b=X_b,
 \qquad v_b\equiv v(\eta)
 \quad\text{on }[-\infty,0]
\]
for every \(b\notin C_v\).

This prescription is a sheaf.  Under restriction to a smaller centre
neighbourhood, only finitely many additional whole branches can be lost.
For gluing, choose one member \(W_0\) of the cover that contains \(\eta\).
The set \(W_0\) contains all but finitely many whole branches, and its
countable exceptional set therefore controls the glued section on all but
countably many whole branches of the union.

Every CLD-smooth function belongs to \(\mathscr C\): the explicit smooth
structure in \cite[\S6.7.3]{CLD25} makes a smooth centre germ constant on
all but finitely many whole branches.  The sheaf \(\mathscr C\) is stable
under decreasing sequential limits.  Indeed, if \(v_n\downarrow v\) on
one open set, the union of the countable exceptional sets of the \(v_n\)
is countable, and outside that union
\[
 v_{n,b}\equiv v_n(\eta)
 \quad\Longrightarrow\quad
 v_b\equiv v(\eta).
\]
This remains true when the last value is \(-\infty\).  The minimality in
\cite[Definition~8.1.1]{CLD25} now gives
\(\operatorname{Psh}_{X_B}\subset\mathscr C\).  Alternatively, the same
argument is a transfinite induction using the common-neighbourhood sequence
at a successor stage and membership in one earlier stage at a limit ordinal
in \cite[Lemma~8.1.4]{CLD25}.  This proves the first assertion and its
ordinal refinement.

We turn to the quantitative assertion.  Put
\[
 c_n=f_n(\eta),
 \qquad d_{n,b}=f'_{n,b}(0).
\]
The topology of \(X_B\) makes \(A\) finite.  The explicit CLD smooth
structure says that the family \((d_{n,b})_{b\in B}\) is finitely supported
and satisfies
\begin{equation}
\label{eq:p071-smooth-centre-balance}
 \sum_{b\in B}d_{n,b}=0.
\end{equation}
Positivity of \(\mathrm d'\mathrm d''f_n\), tested on the calibrated
one-dimensional branches, makes every branch restriction convex.  If
\(b\notin A\), the branch is whole and \(f_{n,b}\) is constant near
\(-\infty\).  Convexity then gives \(d_{n,b}\geq0\) and, for \(r>0\),
\begin{equation}
\label{eq:p071-full-branch-smooth-secant}
 0\leq\frac{c_n-f_{n,b}(-r)}{r}\leq d_{n,b}.
\end{equation}
For \(a\in A\), convexity on \([-\rho_a,0]\) gives
\begin{equation}
\label{eq:p071-exceptional-smooth-secant}
 d_{n,a}\geq
 \frac{c_n-f_{n,a}(-\rho_a)}{\rho_a},
 \qquad
 -d_{n,a}\leq
 \frac{f_{n,a}(-\rho_a)-c_n}{\rho_a}.
\end{equation}
Define
\[
 R_n=\sum_{a\in A}
       \frac{f_{n,a}(-\rho_a)-c_n}{\rho_a}.
\]
For a finite \(J\subset B\setminus A\),
\eqref{eq:p071-smooth-centre-balance},
\eqref{eq:p071-full-branch-smooth-secant}, and
\eqref{eq:p071-exceptional-smooth-secant} yield
\begin{equation}
\label{eq:p071-finite-stage-budget}
 0\leq
 \sum_{b\in J}\frac{c_n-f_{n,b}(-r_b)}{r_b}
 \leq\sum_{b\in J}d_{n,b}
 \leq\sum_{b\notin A}d_{n,b}
 =-\sum_{a\in A}d_{n,a}
 \leq R_n.
\end{equation}
The same calculation with the empty sum on the left shows \(R_n\geq0\).
Since \(f_n\leq f_1\) and \(c_n\geq c\),
\begin{equation}
\label{eq:p071-uniform-exceptional-bound}
 0\leq R_n\leq
 M:=\sum_{a\in A}
       \frac{f_{1,a}(-\rho_a)-c}{\rho_a}<+\infty.
\end{equation}

Suppose now that \(c\in\mathbf R\).  If
\(u_a(-\rho_a)=-\infty\) for one \(a\in A\), its summand in \(R_n\) tends
to \(-\infty\), while every other summand is bounded above by the
corresponding finite summand in \(M\).  This contradicts \(R_n\geq0\).
Thus every exceptional value in \eqref{eq:p071-exceptional-slope-budget}
is finite, \(R_n\to R_A(u;\rho)\), and
\eqref{eq:p071-uniform-exceptional-bound} proves
\eqref{eq:p071-exceptional-slope-budget}.  Passing to the pointwise limit
in \eqref{eq:p071-finite-stage-budget} proves
\eqref{eq:p071-finite-secant-budget}.  In particular, for \(b\notin A\)
and \(r>0\),
\begin{equation}
\label{eq:p071-full-branch-two-sided-bound}
 c-R_A(u;\rho)r\leq u_b(-r)\leq c.
\end{equation}
This proves finiteness on the finite stratum and the first equality in
\eqref{eq:p071-two-endpoint-traces}.

Pointwise decreasing limits preserve convexity and monotonicity wherever
the limit is finite.  Hence \(u_b\) is convex and nondecreasing on the
finite stratum.  Let
\(L_b=\lim_{r\to+\infty}u_b(-r)\), which exists in
\(\mathbf R\cup\{-\infty\}\).  Monotonicity gives
\(u_b(-\infty)\leq L_b\).  On the other hand, \(u_b\leq f_{n,b}\), and
each smooth full-branch restriction is continuous and constant near
\(-\infty\); consequently
\[
 L_b\leq f_{n,b}(-\infty)
 \quad\text{for every }n.
\]
Taking the infimum over \(n\) gives
\(L_b\leq u_b(-\infty)\), and proves the second equality in
\eqref{eq:p071-two-endpoint-traces}.

For a finite convex function on \((-\infty,0]\), the secant quotients in
\eqref{eq:p071-inward-endpoint-slope} increase to the left derivative at
the endpoint as \(r\downarrow0\).  Thus
\eqref{eq:p071-inward-endpoint-slope} follows.  Taking the supremum in
\eqref{eq:p071-finite-secant-budget} independently over the finitely many
radii, and then taking the supremum over finite \(J\), proves
\eqref{eq:p071-summable-endpoint-slopes}.  A summable arbitrarily indexed
family of nonnegative numbers has at most countably many positive entries.
If \(\sigma_b(u)=0\), convexity and monotonicity force \(u_b=c\) on the
finite stratum; the second equality in
\eqref{eq:p071-two-endpoint-traces} then also gives
\(u_b(-\infty)=c\).  This proves the asserted countability of nonconstant
full branches.

Assume next that \(c_n\downarrow-\infty\).  Fix \(b\in B\), let \(I_b\)
be the component of \(V_b\) containing \(0\), and take finite points
\(s<t<0\) in \(I_b\).  Write \(t=(1-\lambda)s+\lambda0\), with
\(0<\lambda<1\).  Convexity gives
\[
 f_{n,b}(t)
 \leq(1-\lambda)f_{n,b}(s)+\lambda c_n
 \leq(1-\lambda)f_{1,b}(s)+\lambda c_n.
\]
The right-hand side tends to \(-\infty\), so \(u_b(t)=-\infty\).  Such an
\(s\) exists for every finite \(t\in I_b\setminus\{0\}\) after moving
slightly farther from the centre inside the open interval \(I_b\).
At the centre the assertion is the hypothesis.  If \(I_b\) also contains
\(-\infty\), then it is the whole branch, and monotonicity gives
\(f_{n,b}(-\infty)\leq f_{n,b}(-r)\leq c_n\) for every \(r>0\), hence
\(u_b(-\infty)=-\infty\).  This proves the stated conclusion on every
component meeting \(\eta\).

It remains to prove sharpness.  Choose a smooth nondecreasing function
\(\psi:[-\infty,0]\to[0,1]\), constant equal to \(0\) near
\(-\infty\), with \(\psi(0)=1\), and put
\[
 \chi(t)=-\int_t^0\psi(s)\,\mathrm ds.
\]
Then \(\chi\leq0\), \(\chi(0)=0\), \(\chi'(0)=1\), and \(\chi\) is smooth,
convex, and constant near \(-\infty\).  Put
\[
 \varepsilon_j=2^{-j},
 \qquad
 S_n=\sum_{j=1}^n\varepsilon_j=1-2^{-n},
 \qquad
 T_n=2^{-n}.
\]
On the centre neighbourhood with \(V_a=(-1,0]\) and every other branch
whole, define
\[
 f_{n,a}(t)=T_n-S_nt,
\]
\[
 f_{n,b_j}(t)=T_n+\varepsilon_j\chi(t)
 \quad(j\leq n),
 \qquad
 f_{n,b}(t)=T_n
 \quad\text{on all remaining branches}.
\]
These functions are branchwise smooth and convex, constant near every
included \(-\infty\), and nonconstant on only finitely many full branches.
Their endpoint derivatives satisfy
\[
 f'_{n,a}(0)+\sum_{j=1}^n f'_{n,b_j}(0)
 =-S_n+\sum_{j=1}^n\varepsilon_j=0.
\]
Thus the explicit CLD smoothness criterion and one-dimensional positivity
give \(f_n\in\operatorname{Psh}_{X_B,0}(V)\).  The sequence decreases.  On
the exceptional branch,
\[
 f_{n+1,a}(t)-f_{n,a}(t)
 =-\varepsilon_{n+1}(1+t)\leq0
 \quad(-1<t\leq0).
\]
On an already active full branch only \(T_n\) decreases; on the newly active
branch \(b_{n+1}\), the inequality follows from \(\chi\leq0\); and it is
immediate on every other branch.  The pointwise limit is
\[
 u_a(t)=-t,
 \qquad
 u_{b_j}(t)=\varepsilon_j\chi(t),
 \qquad
 u_b(t)=0
 \quad\text{otherwise}.
\]
Hence every \(b_j\) is active,
\(\sigma_{b_j}(u)=\varepsilon_j\), and
\[
 \sum_{j\geq1}\sigma_{b_j}(u)=1
 =\frac{u_a(-\rho)-u(\eta)}{\rho}
\]
for every \(0<\rho<1\).  This proves
\eqref{eq:p071-sharp-slope-budget} and completes the proof.
\end{proof}
